\section{全混合代表の六級数と中央平行移動}
preferred中心は前節の$\Psi_{e,m,h}$に固定する。
$\lambda_i^3=1,\Lambda=\prod\lambda_i$に対し
\[
 x_i\mapsto\lambda_i x_i,\quad e_i\mapsto\Lambda^{-1}\lambda_i e_i,\quad
 m_i\mapsto\lambda_i^{-1}m_i,\quad h\mapsto\Lambda h
\]
とし、置換は添字を同時置換して$h$に符号を掛ける。
Weyl対称化はこの有限群作用と可換である。
命題\ref{mix:pure}の各段の原始と曲率を有限群で平均すれば、
係数線形性、有限arity、内部次数、特殊ファイバーの標識を保った共変gaugeを得る。
有理係数性もGalois不変な平均から保たれる。

\begin{theorem}\label{mix:six}
全六非対角一次係数と九二次係数に対して、純曲率代表を
\begin{align*}
 H_h={}&D/3+h^2xyz/3+b(h)\sum_i m_ix_i+\gamma(h)\sum_i e_ix_jx_k\\
 &+U_e(h)\sum_i e_i^2x_i+V_e(h)\sum_i e_je_kx_i
   +R_e(h)\sum_i e_i^3+S_e(h)e_1e_2e_3+T_e(h)\sum_i e_im_i
\end{align*}
と選べる。ここで$b^3=1+h^6/27,b(0)=1$であり、
\begin{gather*}
 \gamma\in1+h^6\mathbb Q[[h^6]],\quad
 U_e,S_e\in h^2\mathbb Q[[h^6]],\quad
 V_e,T_e\in h^4\mathbb Q[[h^6]],\quad R_e\in h^6\mathbb Q[[h^6]],\\
 U_e=-h^2/12+O(h^8),\qquad S_e=h^2/4+O(h^8).
\end{gather*}
六つの未知級数の全値や、この代表の無標識での一意性は主張しない。
\end{theorem}
\begin{proof}
$\Lambda=1$の対角群と巡回群の不変な重み三単項式を列挙すると、
\begin{gather*}
 D,\ xyz,\ \sum e_ix_jx_k,\ \sum e_ix_i^2,\ \sum e_i^2x_i,\\
 \sum e_je_kx_i,\ \sum e_i^3,\ e_1e_2e_3,\ \sum m_ix_i,\ \sum e_im_i
\end{gather*}
の十型となる。
全て置換対称なので係数は$h$の偶級数である。
$D/3$係数を$A$、$\sum e_ix_i^2$係数を$E$とすると
$x_i=A^{-1/3}X_i-(E/A)e_i$で前者を一、後者を零にできる。
$e=0$で三つの独立$m_i$の指標方向はそれぞれ$x_i$のみで、巡回対称性から共通係数となる。
一混合項で定理\ref{mix:b}を用いるとその係数は$b$、純$xyz$は$h^2/3$である。
重みからこれらの係数がさらに$e,m$に依存する余地はない。
残る対角群の$h\mapsto\Lambda h$と置換偶性を合わせると表示の合同条件を得る。
二次初期gaugeを定理\ref{mix:secondcurve}のものに固定すれば、最終線形規格化は
共通拡大$1+O(h^6)$と平行移動$O(h^4)$なので二つの二次係数を変えない。
\end{proof}

\begin{proposition}[全二十一係数と六自由方向]\label{mix:translation}
$x_i$の対角線形項を$L_i$、非対角項を$e_i$、二次項を$m_i$とする。
$X_i=x_i+L_i/2$とし、
\[
 n_i=m_i-L_i^2/4-(e_jL_k+e_kL_j)/2,\qquad
 C_{\rm cl}=\sum_i L_i^3/12+\sum_i e_iL_jL_k/4-\sum_i m_iL_i/2
\]
と置く。共変規格化中心を用いた全混合の代表は
$H_h(X;e,n)+u^3+v^3+C_{\rm cl}$である。
$e=0$の十五係数部分族では追加六級数を用いず
$D(X)/3+h^2X_1X_2X_3/3+b\sum n_iX_i+u^3+v^3+C_{\rm cl}$となる。
全二十二変数源は、$h$を固定する接写像恒等の形式基底変換により、
$L=0$の十六変数源と六次元自由円板との積に分解する。
\end{proposition}
\begin{proof}
中央平行移動は交換子を変えず、実際の中心の共変性により表示の$n,C_{\rm cl}$を与える。
全gaugeは係数を微分せず正規化されているので、追加中央曲率をそのまま残す。
$g=u^3+v^3+C_{\rm cl}$は一意な$M\equiv I_2$により$g\circ M=U^3+V^3$へ移る。
実際$\mathfrak{gl}_2$の軌道微分は
$(\alpha,\beta,\gamma,\delta)\mapsto3\alpha u^3+3\beta u^2v+3\gamma uv^2+3\delta v^3$で可逆である。
混合係数イデアルを$J$とすると$C_{\rm cl}\in J^2$なので$M-I\in J^2$。
従って$(L,e,m,h)\mapsto(L,e\circ M,n\circ M,h)$の接写像は恒等で、形式自己同型となる。
新座標の中心対に$L$は現れず、それを自由座標として残せる。
二座標置換は次数付き、特殊ファイバーで恒等、中心の式そのものを保つので指定周期性も保つ。
PBW定数零など別の中心規格化を用いる場合には、その中央scalar差を同じ向きで二元三次式へ移す必要がある。
\end{proof}

\section{標識付き平面incidenceへの同定}
標的のsquarefree Fermat切片を
\[
 F_{s,r}=F_0+sxyz+\sum_i(P_iu+Q_iv)x_jx_k+\sum_iT_ix_iuv,
 \qquad r=(P,Q,T)
\]
とする。$u=\sum\lambda_ix_i,v=\sum\mu_ix_i$で制限し、三変数marked Hesse規格化した係数を$\sigma$と書く。
\begin{theorem}\label{mix:incidence}
$L=0$の実際の全形式対応は、$h$を固定する$A=\CC[[s,r]]$代数同型により
\[
 B_{16}\simeq A[[\lambda,\mu,h]]/(\sigma(s,r,\lambda,\mu)-h^2/3)
\]
と表示される。全族はこの源に六つの自由変数を加えたものとなる。
原点で$\lambda_i=a_i/3+O(2),\mu_i=c_i/3+O(2)$、$r=(p,q,d)+O(2)$である。
従って平面incidence稿の臨界幾何を実際の形式対応へ適用できる。
\end{theorem}
\begin{proof}
定理\ref{mix:six}の可換代表を$G\equiv I_5$でcanonical Fermat切片へ移す。
元の平面$W_0=(u=v=0)$上の制限は正確に$D/3+h^2xyz/3$である。
新座標の平面は$G^{-1}W_0$であり、切片の制限式はこのHesse式に近恒等三変数線形変換を施したものとなる。
marked Hesse規格化の一意性から$\sigma=h^2/3$を得る。
一次shearは$u\mapsto u-\sum a_ix_i/3,v\mapsto v-\sum c_ix_i/3$なので、表示の接写像となる。
$\partial_s\sigma(0)=1$より右辺は正則十六次元環であり、
$(r,\lambda,\mu,h)$の接写像は可逆である。
形式逆関数定理で同型を得る。全族は命題\ref{mix:translation}から従う。
この同型は形式座標のもので、比較核の有限$h$での代数化を述べていない。
\end{proof}

\begin{corollary}[接階数と正確な臨界イデアル]
分離軸上のcanonical接写像は
$\partial_{p_i}\mapsto\gamma[ux_jx_k]$、
$\partial_{q_i}\mapsto\gamma[vx_jx_k]$、$\partial_{d_i}\mapsto b[x_iuv]$であり、
$L,a,c$の十二方向は零である。従って混合接階数は九、
量子方向も含めた十次minorは$(2h/3)b^3\gamma^6$である。
原点では階数九、分離軸の形式分数体上では十となる。
$L=0$の実際の源における相対臨界イデアルは正確に
$(h,p_1,p_2,p_3,q_1,q_2,q_3)$である。
全族では六つの自由な$L$方向が加わる。
\end{corollary}
\begin{proof}
分離Jacobian環の次数三基底は$xyz,x_iuv,ux_jx_k,vx_jx_k$である。
定理\ref{mix:six}を係数の一次まで取れば接写像とminorを得る。
定理\ref{mix:incidence}で相対臨界イデアルは$(h,\partial_\lambda\sigma,\partial_\mu\sigma)$である。
$h=0$では源の可換代表は古典三次式そのもので、$W_0$上の制限はFermat式、
二つの正規二次類は$\sum p_ix_jx_k,\sum q_ix_jx_k$となる。
内在的臨界条件の証明の完全対合は、これら六係数に対して可逆な線形写像なので、
零スキームのイデアルは$(p,q)$そのものである。
形式座標変更のJacobianは可逆だから、このイデアル同定は集合論だけでなくスキーム論的に成立する。
\end{proof}
